\subsection{Fourier 积分的收敛与可和性}
\label{sec:ch1-fourier-integral-convergence-summability}

从 Fourier 变换恢复函数的问题, 与 Fourier 级数中的同一问题相似. 需要确定在何时以及是否有
\[
 \lim_{R\to\infty}\int_{B_R}\widehat f(\xi)e^{2\pi ix\cdot\xi}\,d\xi=f(x),
\]
其中 $B_R=\{Rx:x\in B\}$, $B$ 是原点的一个开凸邻域, 而极限或者按 $L^p$ 意义理解, 或按几乎处处逐点意义理解. 若定义部分和算子 $S_R$ 为
\[
 \widehat{(S_Rf)}=\chi_{B_R}\widehat f,
\]
则这个问题等价于确定是否有
\[
 \lim_{R\to\infty}S_Rf=f.
\]
与引理~\ref{lem:ch1-lp-partial-sum-convergence} 类似, 范数收敛的充要条件是
\[
 \|S_Rf\|_p\le C_p\|f\|_p,
\]
其中 $C_p$ 与 $R$ 无关. 当 $n=1$ 时这一点成立; 我们将在 \MBUnresolvedReference{chapter}{第 3 章} 证明. 当 $n>1$ 时, 还将证明几个部分结果, 但一般而言, 当 $p\ne2$ 时不存在范数收敛. 这一点将在 \MBUnresolvedReference{chapter}{第 8 章} 讨论.

当 $n=1$ 且 $B=(-1,1)$ 时,
\[
 S_Rf(x)=D_R*f(x),
\]
\hypertarget{chapter-01-page-030}{}
其中 $D_R$ 是 Dirichlet 核,
\[
 D_R(x)=\int_{-R}^{R}e^{2\pi ix\xi}\,d\xi
 =\frac{\sin(2\pi Rx)}{\pi x}.
\]
它显然不可积, 但对任意 $q>1$ 都属于 $L^q(\mathbb R)$, 所以当 $f\in L^p$, $1<p<\infty$ 时, $D_R*f$ 有定义.

几乎处处收敛取决于估计
\[
 \left\|\sup_R|S_Rf|\right\|_p\le C_p\|f\|_p.
\]
当 $1<p<\infty$ 时此式成立(Carleson--Hunt 定理), 但这里无法证明.

对 Fourier 变换, Cesàro 可和法是取部分和算子的积分平均
\[
 \sigma_Rf(x)=\frac1R\int_0^RS_tf(x)\,dt,
\]
并确定是否有 $\lim\sigma_Rf(x)=f(x)$. 当 $n=1$ 且 $B=(-1,1)$ 时,
\[
 \sigma_Rf(x)=F_R*f(x),
\]
其中 $F_R$ 是 Fejér 核,
\begin{equation}
\label{eq:ch1-fourier-integral-fejer-kernel}
 F_R(x)=\frac1R\int_0^RD_t(x)\,dt
 =\frac{\sin^2(\pi Rx)}{R(\pi x)^2}.
\end{equation}
与 Dirichlet 核不同, Fejér 核可积. 由于它具有与 \eqref{eq:ch1-fejer-kernel-properties} 类似的性质, 可以证明在 $L^p$, $1\le p<\infty$ 中,
\[
 \lim_{R\to\infty}\sigma_Rf=f.
\]
证明与定理~\ref{thm:ch1-fejer-norm-convergence} 的证明类似. 在第~\ref{chap:hardy-littlewood-maximal-function}~章, 我们将证明两个一般结果, 并由此推出这种可和法以及下面的可和法在 $L^p$ 中和几乎处处逐点意义下的收敛性.

Abel--Poisson 可和法是在反演公式中引入因子 $e^{-2\pi t|\xi|}$. 对每个 $t>0$, 该积分收敛, 然后令 $t$ 趋于 $0$. 若改为引入因子 $e^{-\pi t^2|\xi|^2}$, 就得到 Gauss--Weierstrass 可和法. 更准确地, 定义
\begin{equation}
\label{eq:ch1-abel-poisson-summability}
 u(x,t)=\int_{\mathbb R^n}e^{-2\pi t|\xi|}\widehat f(\xi)e^{2\pi ix\cdot\xi}\,d\xi,
\end{equation}
\begin{equation}
\label{eq:ch1-gauss-weierstrass-summability}
 w(x,t)=\int_{\mathbb R^n}e^{-\pi t^2|\xi|^2}\widehat f(\xi)e^{2\pi ix\cdot\xi}\,d\xi,
\end{equation}
并试图确定是否有
\begin{equation}
\label{eq:ch1-abel-poisson-limit}
 \lim_{t\to0^+}u(x,t)=f(x),
\end{equation}
\hypertarget{chapter-01-page-031}{}
\begin{equation}
\label{eq:ch1-gauss-weierstrass-limit}
 \lim_{t\to0^+}w(x,t)=f(x),
\end{equation}
其中极限在 $L^p$ 中取, 或几乎处处逐点取.

可以证明 $u(x,t)$ 在半空间 $\mathbb R_+^{n+1}=\mathbb R^n\times(0,\infty)$ 中调和. 当 $n=1$ 时, 有一个与 \eqref{eq:ch1-fourier-series-unit-disk} 类似的等价公式:
\begin{equation}
\label{eq:ch1-abel-poisson-complex-formula}
 u(z)=\int_0^\infty\widehat f(\xi)e^{2\pi iz\xi}\,d\xi
 +\int_{-\infty}^0\widehat f(\xi)e^{2\pi i\overline z\xi}\,d\xi,
 \qquad z=x+it,
\end{equation}
它立即蕴含 $u$ 调和. 极限 \eqref{eq:ch1-abel-poisson-limit} 可解释为 Dirichlet 问题的边界条件
\[
 \Delta u=0\quad\text{on }\mathbb R_+^{n+1},
 \qquad
 u(x,0)=f(x),\quad x\in\mathbb R^n.
\]

由 \eqref{eq:ch1-abel-poisson-summability},
\[
 u(x,t)=P_t*f(x),
\]
其中 $\widehat P_t(\xi)=e^{-2\pi t|\xi|}$. 当 $n=1$ 时可通过简单计算证明, 当 $n>1$ 时可通过更困难的计算证明(见 Stein 与 Weiss \MBUnresolvedReference{citation}{[18, p. 6]}),
\begin{equation}
\label{eq:ch1-poisson-kernel-half-space}
 P_t(x)=\frac{\Gamma\left(\frac{n+1}{2}\right)}{\pi^{(n+1)/2}}
 \frac{t}{(t^2+|x|^2)^{(n+1)/2}}.
\end{equation}
这称为 Poisson 核.

在 Gauss--Weierstrass 可和法的情形, 可以证明函数 $\widetilde w(x,t)=w(x,\sqrt{4\pi t})$ 是热方程
\[
 \frac{\partial\widetilde w}{\partial t}-\Delta\widetilde w=0
 \quad\text{on }\mathbb R_+^{n+1},
 \qquad
 \widetilde w(x,0)=f(x),\quad x\in\mathbb R^n
\]
的解, 而 \eqref{eq:ch1-gauss-weierstrass-limit} 可解释为该问题的初值条件. 还有公式
\[
 w(x,t)=W_t*f(x),
\]
其中 $W_t$ 是 Gauss--Weierstrass 核,
\begin{equation}
\label{eq:ch1-gauss-weierstrass-kernel}
 W_t(x)=t^{-n}e^{-\pi|x|^2/t^2}.
\end{equation}
此公式可由引理~\ref{lem:ch1-gaussian-fourier-transform} 和 \eqref{eq:ch1-fourier-dilation} 证明.

